% Generated draft: requires mathematical and source review.
\documentclass[11pt]{article}
\usepackage[margin=25mm]{geometry}
\usepackage{fontspec}
\setmainfont{Latin Modern Roman}
\usepackage{amsmath,amssymb,mathtools,unicode-math}
\setmathfont{Latin Modern Math}
\usepackage{enumitem,needspace,xurl,hyperref}
\setlength{\parindent}{0pt}
\setlength{\parskip}{5pt}
\setlength{\emergencystretch}{3em}
\newcommand{\N}{\mathbb{N}}
\newcommand{\Z}{\mathbb{Z}}
\newcommand{\Q}{\mathbb{Q}}
\newcommand{\R}{\mathbb{R}}
\newcommand{\C}{\mathbb{C}}
\newcommand{\F}{\mathbb{F}}
\newcommand{\A}{\mathbb{A}}
\newcommand{\PP}{\mathbb{P}}
\newcommand{\E}{\mathbb{E}}
\newcommand{\eps}{\varepsilon}
\newcommand{\dd}{\,\mathrm{d}}
\begin{document}
\section*{UnsolvedMath ID: 3100062}
\textbf{Source problem code:} catalogue entry
\subsection{Definitions and mathematical statement}
Write \(\mathbb Z_{>0}=\{1,2,\ldots\}\). Under the usual integer-coefficient interpretation of the source, a linear homogeneous equation in \(n\) variables has the form
\[
 a_1x_1+\cdots+a_nx_n=0,\qquad a=(a_1,\ldots,a_n)\in\mathbb Z^n.
\]
For a positive integer \(r\), call this equation \(r\)-Ramsey on the positive integers if, for every coloring \(c:\mathbb Z_{>0}\to\{1,\ldots,r\}\), there are \(x_1,\ldots,x_n\in\mathbb Z_{>0}\) satisfying the equation and \(c(x_1)=\cdots=c(x_n)\). The variables are allowed to coincide in this conventional partition-regularity definition. Call the equation partition regular if it is \(r\)-Ramsey for every positive integer \(r\).

Rado's question asks whether
\[
\forall n\ge1\ \exists M(n)\ge1\ \forall a\in\mathbb Z^n:
\quad\bigl(a\text{ is }M(n)\text{-Ramsey}\bigr)
\Longrightarrow
\bigl(\forall r\ge1,\ a\text{ is }r\text{-Ramsey}\bigr).
\]
Here \(M(n)\) must depend only on the number of variables, not on the coefficients. Allowing the identically zero equation does not affect the question, since it is already partition regular. Rational coefficients give the same class after clearing denominators. The source does not explicitly specify the coefficient domain or a distinctness convention; the displayed formulation records the standard integer-coefficient, repetitions-allowed reading, which requires source confirmation if another reading was intended.

The supplied summary mentions Rado's characterization of partition-regular linear equations, but does not claim that it settles this uniform cutoff question. It marks the specific question open/unverified, based on no verified resolution. The source page's open label is from its October 2020 snapshot and was not checked item by item against later literature.
\subsection{Short English statement}
Is there a number of colors depending only on the number of variables such that an equation which always has a monochromatic positive-integer solution for that many colors automatically has one for every finite number of colors? The usual integer-coefficient interpretation is made explicit, and the source's unverified open status is retained.
\subsection{Sources}
\begin{itemize}
\item Local registry: \url{data/unsolvedmath/problems.json}, permanent UnsolvedMath ID \(3100062\), source code catalogue entry, attributed to Rado; supplied statement, background, and research summary, updated 2026-08-06.
\item Cooper, \emph{Combinatorial Problems I Like}, October 2020 snapshot, Combinatorial Number Theory, source-order bullet 62: \url{https://people.math.sc.edu/cooper/combprob.html}. HTML extraction accessed 2026-07-29.
\item The supplied research summary refers to Rado's theorem on partition regularity but supplies no original-paper citation or verified resolution of the cutoff question.
\end{itemize}
\end{document}
